\subsubsection*{2.4.6.（原文第二次出现）}
设 $X_{s(x)}$ 为 $X_s$ 在 $x$ 处的 Hensel 化, 置
\[
X^*_{s(x)}=X_{s(x)}-\{x\}.
\]
$(X_{s(x)},\{x\})$ 分别以 $K_s$, $R\Psi_{\bar\eta}(K)$ 和 $R\Phi_{\bar\eta}(K)$ 为系数的长正合超上同调列, 给出如下长正合列的交换图:
\begin{equation}\tag{2.4.6.1}
\resizebox{\textwidth}{!}{%
\begin{tikzcd}[ampersand replacement=\&, column sep=large, row sep=large]
{}
  \& 0 \arrow[r] \arrow[d]
  \& R^{n-1}\Phi_{\bar\eta}(K)_x \arrow[r] \arrow[d]
  \& R^{n-1}\Phi_{\bar\eta}(K)_x \arrow[r] \arrow[d]
  \& 0 \arrow[d]
  \& {} \\
{} \arrow[r]
  \& \mathbb H^{n-1}(X^*_{s(x)},K)
     \arrow[r, "j"] \arrow[d, equal]
  \& \mathbb H^n_{\{x\}}(X_s,K) \arrow[r] \arrow[d]
  \& \underline{H}^{n}(K)_x \arrow[r] \arrow[d]
  \& \mathbb H^n(X^*_{s(x)},K)
     \arrow[r] \arrow[d, equal]
  \& {} \\
{} \arrow[r]
  \& \mathbb H^{n-1}(X^*_{s(x)},K)
     \arrow[r, "j", "{\mbox{\textcircled{1}}}"'] \arrow[d]
  \& \mathbb H^n_{\{x\}}\bigl(X_s,R\Psi_{\bar\eta}(K)\bigr)
     \arrow[r] \arrow[d]
  \& R^n\Psi_{\bar\eta}(K)_x \arrow[r] \arrow[d]
  \& \mathbb H^n(X_{s(x)},K) \arrow[r] \arrow[d]
  \& {} \\
{}
  \& 0 \arrow[r]
  \& R^n\Phi_{\bar\eta}(K)_x \arrow[r, equal]
  \& R^n\Phi_{\bar\eta}(K)_x \arrow[r]
  \& 0
  \& {}
\end{tikzcd}}
\end{equation}
箭头 \textcircled{1} 的像包含于惯性不变部分中.

原文将本段和前一段都编号为 2.4.6, 也将两个主要图都编号为 (2.4.6.1). 这里保留重复的编号; 第 XV 讲 2.2.6 所引用的是第二次出现的这一处.

\renewcommand{\refname}{参考文献}
\begin{thebibliography}{9}
\bibitem{GiraudClassifyingTopos}
J.~Giraud, ``Classifying topos,'' pp.~43--56 in
\emph{Toposes, Algebraic Geometry and Logic},
Lecture Notes in Mathematics 274, Springer-Verlag, 1971.

\bibitem{NagataImbedding}
M.~Nagata, ``A generalization of the imbedding problem,''
\emph{J. Math. Kyoto} \textbf{3}, no.~1 (1963), pp.~89--102.
\end{thebibliography}
